\chapter*{Úvod}
\addcontentsline{toc}{chapter}{Úvod}

V úvodu závěrečné práce autor vysvětlí, proč si vybral zvolené téma, tedy \textbf{motivaci} celé závěrečné práce. V úvodu nesmí chybět přesně formulovaný \textbf{hlavní cíl} závěrečné práce (popř. dílčí cíle), měla by zde být nastíněna \textbf{metodika} celé závěrečné práce (popř. výzkumné otázky či hypotézy). Zvykem bývá rovněž nastínit \textbf{hlavní výsledky/výstupy} závěrečné práce. 

% TODO nahoře
% TODO zdroje v textu

Digitální nástroje určené pro řízení pracovního prostředí (workplace management) se v posledních letech staly nedílnou součástí moderních organizací. S nástupem hybridního pracovního modelu, rostoucí flexibilitou výkonu práce a proměnou způsobu využívání kancelářských prostor se zvyšují nároky na software, který zaměstnancům umožní efektivně plánovat přítomnost v kanceláři, rezervovat pracovní místa či řešit administrativní povinnosti spojené s výkonem práce na dálku (Deskbird, b.r.; UC Today, 2026). Kromě funkční správnosti je však stále větší důraz kladen na přehlednost, intuitivnost a uživatelskou přívětivost digitálních řešení.

Předmětem této práce je existující aplikace pro správu pracovního prostředí, která vznikala postupně jako interní nástroj podporující provozní procesy organizace, jejíž vývoj a nasazení bylo urychleno nástupem pandemie COVID-19. V tomto období představovala aplikace rychle navržený interní nástroj se základními funkcemi, jehož primárním cílem bylo umožnit zaměstnancům rezervovat pracovní místa a evidovat práci z domova. Její vývoj se proto historicky soustředil primárně na okamžitou funkčnost a podporu provozních procesů organizace, nikoliv na vizuální podobu či konzistentní uživatelskou zkušenost. Postupem času však byla aplikace průběžně rozšiřována o další moduly, které reflektovaly nové interní i legislativní požadavky organizace. Mezi nejvýznamnější doplnění patřila například správa zasedacích místností, evidence služebních cest, přesčasů, pohotovostí či podpora plánování směn pro pracovníky zaměstnané formou DPP/DPČ.

V důsledku tohoto organického růstu dnes aplikace obsahuje široké spektrum funkcionalit, které již dalece přesahují její původní účel. S rostoucím rozsahem se však výrazně zvýšila i komplexita jejího uživatelského rozhraní. Jednotlivé moduly byly vytvářeny v různých časových obdobích a často různými autory, přičemž chyběl jednotný designový rámec či definovaný design systém. Rozhraní aplikace se tak postupně stalo nekoherentním, vizuálně nejednotným a pro uživatele obtížněji srozumitelným. Mnohé prvky postrádají logické uspořádání, vizuální hierarchii nebo intuitivně pochopitelné ovládání. Tyto nedostatky se promítají do obtížné orientace uživatelů, zvýšené kognitivní zátěže a neefektivních pracovních postupů. S ohledem na klíčovou roli aplikace v interních procesech proto vyvstává reálná potřeba zlepšit její použitelnost, vizuální konzistenci a celkovou přehlednost tak, aby lépe podporovala každodenní pracovní úkony.

Tuto potřebu umocňuje i existující strategický záměr transformovat aplikaci z interního nástroje do plnohodnotného komerčního produktu. Komerční využití s sebou nese nutnost poskytovat profesionální a použitelný produkt, který bude schopen oslovit širší spektrum zákazníků a konkurovat zavedeným řešením na trhu workplace managementu.

Cílem předložené diplomové práce je proto navrhnout a ověřit nový uživatelský design, který zvýší přehlednost, intuitivnost a vizuální konzistenci rozhraní, a současně umožní efektivní budoucí škálování systému pro účely zmíněné komercializace. Redesign představuje komplexní proces zahrnující analýzu současného stavu, vytvoření sjednocené vizuální identity, návrh uceleného design systému a revizi architektury. V rámci této práce je návrh realizován do fáze uživatelsky otestovaných interaktivních prototypů klíčových obrazovek, které jsou metodicky i vizuálně připraveny k následné technické implementaci vývojovým týmem.

Pro naplnění tohoto cíle byly stanoveny následující výzkumné otázky:

\begin{itemize}
    \item \textbf{VO1:} Jaké jsou hlavní bariéry v použitelnosti a orientaci v současném uživatelském rozhraní aplikace z pohledu interních uživatelů?
    \item \textbf{VO2:} Jak navrhnout design systém, aby zajistil vizuální konzistenci napříč moduly a umožnil efektivní budoucí rozšiřování aplikace o nové funkcionality?
    \item \textbf{VO3:} V jaké míře snižuje redesignované rozhraní kognitivní zátěž a zvyšuje intuitivnost používání oproti původnímu stavu?
\end{itemize}

Text práce je logicky strukturován do X kapitol. Po úvodním představení teoretických východisek uživatelsky orientovaného designu následuje analytická část věnovaná uživatelskému výzkumu. Praktická část následně popisuje tvorbu design systému a interaktivních prototypů klíčových obrazovek, jejichž úspěšnost je v závěru práce evaluována prostřednictvím reportu z testování použitelnosti.